% UBT-AI-PROVENANCE-BEGIN
% schema: ubt-ai-provenance/v1
% tier: C_working
% ai_assistance: disclosed
% human_review: risk-based
% editorial_responsibility: Ing. David Jaroš
% policy: ../../../AI_PROVENANCE.md
% notice: Working material; exhaustive human review is not claimed.
% UBT-AI-PROVENANCE-END

\chapter{Pokračování do komplexního času a Wardovy identity}
\label{ch:ct-scheme}

\section{Proč tato kapitola zůstává v učebnici}
Starší práce na konstantě alfa používaly „schéma CT“ založené na analytickém
pokračování \(\tau=t+i\psi\). Archivovaná implementace není aktuálním zdrojem
tvrzení, avšak základní matematika --- analytické pokračování, kalibrační
identity a limita reálného času --- je standardní a užitečná. Tato kapitola
proto nejprve vysvětluje jednotlivé prvky a poté přesně vymezuje, co pro UBT
zakládají a co nikoli.

\section{Komplexní čas jako parametr}
Pro analytickou veličinu $F(\tau)$ umožňuje vztah
\[
 \tau=t+i\psi
\]
přecházet mezi oscilačním a tlumeným popisem. Například pro energetický vlastní
mód platí
\[
 e^{-iE\tau/\hbar}=e^{-iEt/\hbar}e^{+E\psi/\hbar},
\]
přičemž znaménko tlumení určuje konvence imaginárního směru. V konvenci theta
použité v dalších kapitolách je
\[
 e^{-\pi(\psi-it)n^2}
 =e^{-\pi\psi n^2}e^{i\pi t n^2}.
\]
Žádná z těchto analytických identit sama o sobě nevytváří novou renormalizační
konstantu ani neurčuje fyzikální vazbu.

\section{Wardova--Takahashiho identita: standardní argument}
Nechť je v QED $S(p)$ úplným fermionovým propagátorem a
$\Gamma^\mu(p+q,p)$ úplným elektromagnetickým vrcholem. Kalibrační invariance
dává Wardovu--Takahashiho identitu
\[
 q_\mu\Gamma^\mu(p+q,p)
 =S^{-1}(p+q)-S^{-1}(p).
\]
V limitě $q\to0$ dostáváme
\[
 \Gamma^\mu(p,p)=\frac{\partial S^{-1}(p)}{\partial p_\mu}.
\]
Na úrovni renormalizačních konstant z toho plyne
\[
 \boxed{Z_1=Z_2,}
\]
za předpokladu, že regulátor a odečítací předpis zachovávají kalibrační
symetrii. Jde o standardní výsledek kalibrační teorie. Na pokračování do
komplexního času jej lze použít teprve po důkazu, že zvolená kontura a regulátor
zachovávají předpoklady této identity.

\section{Transverzalita vlastní energie fotonu}
Ze stejné kalibrační symetrie plyne
\[
 q_\mu\Pi^{\mu\nu}(q)=0,
\]
takže vakuová polarizace má až na volbu konvencí tvar
\[
 \Pi^{\mu\nu}(q)
 =\left(q^2g^{\mu\nu}-q^\mu q^\nu\right)\Pi(q^2).
\]
Renormalizaci náboje lze proto vyjádřit pomocí příčné dvoubodové funkce fotonu.
Opět jde o standardní omezení kvantové teorie pole, nikoli o teorém specifický
pro UBT.

\section{Kontrolované propojení s UBT}
Konstrukce renormalizace UBT v komplexním čase by musela na základě skutečné
akce a skutečné integrační kontury doložit:
\begin{enumerate}
 \item dostatečnou analyticitu pro pokračování mezi zvolenou konturou v $\psi$
       a řezem reálného času;
 \item zachování příslušných kalibračních nebo BRST identit;
 \item regularizační a protitermové schéma slučitelné s těmito identitami;
 \item spojitý přechod ke standardní QED v příslušné limitě reálného času;
 \item výslovný výpočet každého konečného zbytku specifického pro schéma CT,
       nikoli jeho určení definicí.
\end{enumerate}
Teprve poté by bylo možné povýšit faktor specifický pro CT z modelového
předpokladu na odvozenou veličinu.

\section{Proč se archivní tvrzení \(\mathcal R_{\rm UBT}=1\) nepřebírá}
Historická větev programu alfa uváděla za předpokladů A1--A3 dvousmyčkovou
výchozí hodnotu \(\mathcal R_{\rm UBT}=1\) a následně označila výslednou základní
hodnotu alfa za „bez přizpůsobování parametrů“. Současné stavové dokumenty úrovně A již
nedovolují vydávat toto znění za dokončené odvození konstanty alfa: propojení
koeficientu v mezeře G137-B zůstává otevřené a audit efektivní násobnosti není
uzavřen. Aktuální učebnice proto zachovává matematiku Wardových identit, ale
archivované tvrzení znovu nepoužívá jako platný zdroj současného stavu.

Jde o užitečný příklad obecného pravidla: správná podmíněná identita kvantové
teorie pole může zůstat platná, i když je širší fenomenologická interpretace,
která na ní byla vystavěna, přeřazena na nižší úroveň jistoty.

\section{Vazba na kapitolu o theta funkcích}
Redukované jádro theta
\[
 S(t,\psi)=\sum_n a_n e^{-\pi\psi n^2}e^{i\pi t n^2}
\]
poskytuje čistší kontrolovanou laboratoř pro komplexní čas než starší
renormalizační cesta programu alfa. Parametr $\psi$ zde působí jako proměnná
tepelného času, zatímco $t$ nese oscilační fázi. Fourierovu, Laplaceovu,
Mellinovu a Poissonovu--Jacobiho transformaci lze poté přesně odvodit ještě
předtím, než se k nim připojí jakékoli fyzikální tvrzení UBT.
